% !TEX TS-program = pdflatex
% !TEX encoding = UTF-8 Unicode

% This file is a template using the "beamer" package to create slides for a talk or presentation
% - Giving a talk on some subject.
% - The talk is between 15min and 45min long.
% - Style is ornate.

% MODIFIED by Jonathan Kew, 2008-07-06
% The header comments and encoding in this file were modified for inclusion with TeXworks.
% The content is otherwise unchanged from the original distributed with the beamer package.

\documentclass{beamer}
\setbeamercovered{invisible}



% Copyright 2004 by Till Tantau <tantau@users.sourceforge.net>.
%
% In principle, this file can be redistributed and/or modified under
% the terms of the GNU Public License, version 2.
%
% However, this file is supposed to be a template to be modified
% for your own needs. For this reason, if you use this file as a
% template and not specifically distribute it as part of a another
% package/program, I grant the extra permission to freely copy and
% modify this file as you see fit and even to delete this copyright
% notice. 


\mode<presentation>
{
  \usetheme{Warsaw}
  % or ...

  % or whatever (possibly just delete it)
}


\usepackage[english]{babel}
% or whatever

\usepackage[utf8]{inputenc}
% or whatever

\usepackage{times}
\usepackage[T1]{fontenc}
% Or whatever. Note that the encoding and the font should match. If T1
% does not look nice, try deleting the line with the fontenc.

%%% MATH RELATED


%%% AMS math stuff
\usepackage{amsmath}
\usepackage{amssymb}
\usepackage{amsthm}
\usepackage{mathrsfs}
\usepackage{enumerate}

%%% Theorem environments
\newtheorem{thm}{Theorem}[subsection]
\newtheorem{prop}[thm]{Proposition}
\newtheorem{cor}{Corollary}[thm]
\newtheorem{por}[thm]{Porism}
\newtheorem{lem}[thm]{Lemma}
\theoremstyle{definition}
\newtheorem{prob}[thm]{Problem}
\newtheorem{soln}{Solution}
\newtheorem{defn}[thm]{Definition}
\newtheorem{ex}[thm]{Example}
\newtheorem{quest}[thm]{Question}
\newtheorem{remk}[thm]{Remark}

%%% Typesetting shortcuts
\newcommand{\tn}[1]{\textnormal{#1}}
\newcommand{\ol}[1]{\overline{#1}}
\newcommand{\wt}[1]{\widetilde{#1}}
\newcommand{\wh}[1]{\widehat{#1}}
\newcommand{\vocab}[1]{\textbf{#1}\index{#1}}

%%% Math shortcuts
\newcommand{\bbr}{\mathbb R}
\newcommand{\bbz}{\mathbb Z}
\newcommand{\bbq}{\mathbb Q}
\newcommand{\bbn}{\mathbb N}
\newcommand{\bbf}{\mathbb F}
\newcommand{\bbc}{\mathbb C}
\newcommand{\bbd}{\mathbb D}
\newcommand{\bba}{\mathbb A}
\newcommand{\bbp}{\mathbb P}
\newcommand{\bbg}{\mathbb G}
\newcommand{\bbv}{\mathbb V}
\newcommand{\dih}[1]{\mathcal D_{#1}}
\newcommand{\sym}[1]{\mathcal S_{#1}}
\newcommand{\vspan}{\tn{span}}
\newcommand{\trace}{\tn{trace}}
\newcommand{\diff}{\backslash}
\newcommand{\stab}{\tn{stab}}
\newcommand{\conv}{\tn{conv}}
\newcommand{\img}{\tn{img}}
\newcommand{\coker}{\tn{coker}}
\newcommand{\id}{\tn{id}}
\newcommand{\Hom}{\tn{Hom}}
\newcommand{\End}{\tn{End}}
\newcommand{\Aut}{\tn{Aut}}
\newcommand{\aut}{\tn{Aut}}
\newcommand{\ann}{\tn{Ann}}
\newcommand{\GL}{\tn{GL}}
\newcommand{\Gr}{\tn{Gr}}
\newcommand{\lord}{\preccurlyeq}
\newcommand{\rord}{\succcurlyeq}
\newcommand{\tr}{\textnormal{Tr}}
\newcommand{\Tr}{\textnormal{Tr}}
\newcommand{\bbl}{\mathbb{L}}
\newcommand{\C}{\mathscr{C}}
\newcommand{\X}{\mathscr{X}}
\renewcommand{\S}{\mathscr{S}}
\newcommand{\M}{\mathscr{M}}
\renewcommand{\L}{\mathcal{L}}


%%% Algebraic Geometry
\newcommand{\height}{\textnormal{ht}}
\newcommand{\A}{\mathbb{A}}
\newcommand{\p}{\mathfrak{p}}
\newcommand{\sheaf}[1]{\mathcal{#1}}
\newcommand{\spec}{\textnormal{Spec}}
\newcommand{\proj}{\textnormal{Proj}}
\newcommand{\Aff}{\textnormal{Aff}}
\newcommand{\skel}{\textnormal{skel}}
\newcommand{\supp}{\textnormal{supp}}
\newcommand{\orb}{\textnormal{orb}}
\newcommand{\Proj}{\textnormal{Proj}}
\newcommand{\Pic}{\textnormal{Pic}}
\newcommand{\Rees}{\textnormal{Rees}}
\newcommand{\shom}{\mathcal{H}om}

\renewcommand*\arraystretch{1.3}

%%% paper specific definitions
\newcommand{\weyl}{\Omega}
\newcommand{\weyll}{{\widehat{\Omega}}}
\newcommand{\weylll}{{\widetilde{\Omega}}}
\newcommand{\mweyl}{{M_N(\Omega)}}
\newcommand{\mweyll}{{M_N(\widehat{\Omega})}}
\newcommand{\mweylll}{{M_N(\widetilde{\Omega})}}
\newcommand{\seq}{\text{Seq}}
\newcommand{\tail}{\text{Tail}}
\newcommand{\Ad}{\textnormal{Ad}}
\newcommand{\sech}{\textnormal{sech}}
\newcommand{\colim}{\varinjlim}
\newcommand{\limit}{\varprojlim}
\newcommand{\Bis}{\textnormal{Bis}}
\newcommand{\m}{\mathfrak{m}}
\newcommand{\mxx}[4]{\left(\begin{array}{cc} #1 & #2\\ #3 & #4 \end{array}\right)}
\newcommand{\diag}{\text{diag}}
\newcommand{\qdet}{\textnormal{qdet}}
\newcommand{\mdet}{\textnormal{mdet}}
\newcommand{\mtau}{\mathcal{T}}
\newcommand{\cof}{\textnormal{cof}}
\newcommand{\minor}{\textnormal{minor}}
\newcommand{\holo}{Holo}
\newcommand{\ord}{\textnormal{order}}
\newcommand{\mult}{\mathfrak M}






\title{MATH 350-2 Advanced Calculus} 
\subtitle
{} % (optional)

\author[W.R. Casper] % (optional, use only with lots of authors)
{W.R. Casper}
% - Use the \inst{?} command only if the authors have different
%   affiliation.

\institute[California State University Fullerton] % (optional, but mostly needed)
{
  Department of Mathematics\\
  California State University Fullerton}
% - Use the \inst command only if there are several affiliations.
% - Keep it simple, no one is interested in your street address.

\subject{Talks}
% This is only inserted into the PDF information catalog. Can be left
% out. 



% If you have a file called "university-logo-filename.xxx", where xxx
% is a graphic format that can be processed by latex or pdflatex,
% resp., then you can add a logo as follows:

% \pgfdeclareimage[height=0.5cm]{university-logo}{university-logo-filename}
% \logo{\pgfuseimage{university-logo}}



% Delete this, if you do not want the table of contents to pop up at
% the beginning of each subsection:
\AtBeginSubsection[]
{
  \begin{frame}<beamer>{Outline}
    \tableofcontents[currentsection,currentsubsection]
  \end{frame}
}


% If you wish to uncover everything in a step-wise fashion, uncomment
% the following command: 

%\beamerdefaultoverlayspecification{<+->}


\begin{document}

\begin{frame}
  \titlepage
\end{frame}

\begin{frame}{Outline}
  \tableofcontents
  % You might wish to add the option [pausesections]
\end{frame}

% Since this a solution template for a generic talk, very little can
% be said about how it should be structured. However, the talk length
% of between 15min and 45min and the theme suggest that you stick to
% the following rules:  

% - Exactly two or three sections (other than the summary).
% - At *most* three subsections per section.
% - Talk about 30s to 2min per frame. So there should be between about
%   15 and 30 frames, all told.

\section{Real Analysis Lecture 21}

\subsection{Series tests}

\begin{frame}{Tests for Convergence}
\begin{thm}[Comparison Test]
If $a_n>0$ and $b_n > 0$ and there exists $c>0$ and $N\in\mathbb Z_+$ with $a_n < cb_n$ for all $n\geq N$ then
$$\sum b_n\ \ \text{converges}\quad \Rightarrow\quad\sum a_n\ \ \text{converges}$$
\end{thm}
\end{frame}

\begin{frame}{Tests for Convergence}
\begin{thm}[Limit Comparison Test]
If $a_n>0$ and $b_n > 0$ and suppose that $\lim a_n/b_n = c > 0$.  Then
$$\sum b_n\ \ \text{converges}\quad \Leftrightarrow\quad\sum a_n\ \ \text{converges}$$
\end{thm}
\end{frame}

\begin{frame}{Limit infimum/limit supremum}
\begin{defn}
Let $a_n$ be a sequence.  Then
$$\liminf a_n = \lim_n \inf\{a_k: k\geq n\}$$
$$\limsup a_n = \lim_n \sup\{a_k: k\geq n\}$$
\end{defn}
\begin{thm}
If $a_n$ is bounded, then $\limsup a_n$ and $\liminf a_n$ exist.
\end{thm}
More generally, we write $\liminf a_n = \pm\infty$ and $\limsup a_n = \pm \infty$.
\end{frame}

\begin{frame}{Tests for Convergence}
\begin{thm}[Ratio Test]
Let $r = \liminf |a_{n+1}/a_n|$ and $R = \limsup |a_{n+1}/a_n|$.
\begin{enumerate}[(a)]
\item if $R < 1$ then the series is absolutely convergent
\item if $r > 1$ then the series is divergent
\item if $r \leq 1\leq R$ the test is inconclusive
\end{enumerate}
\end{thm}
\end{frame}

\begin{frame}{Tests for Convergence}
\begin{thm}[Root Test]
Let $\rho = \limsup \sqrt[n]{|a_n|}$.
\begin{enumerate}[(a)]
\item if $\rho < 1$ then the series is absolutely convergent
\item if $\rho > 1$ then the series is divergent
\item if $\rho = 1$ the test is inconclusive
\end{enumerate}
\end{thm}
\end{frame}

\subsection{Infinite products}

\begin{frame}{Basic definition}
\begin{defn}
An \textbf{infinite product} is an ordererd pair of sequences $(\{a_n\},\{p_n\})$ where 
$$p_n = a_1a_2\dots a_n.$$
The sequence $\{p_n\}$ is called the \textbf{sequence of partial products} and $p_n$ is called the $n$'th partial product.
We say the series converges or diverges if $\{p_n\}$ converges or diverges.
\end{defn}
\textbf{Notation:}
$$a_1a_2a_3\dots\quad\text{or}\quad \prod_{n=1}^\infty a_n.$$
\end{frame}

\begin{frame}{Basic definition}
We write
$$\prod_{n=1}^\infty a_n = L$$
to mean that the limit of the sequence $p_n$ is $L$.
We call $L$ the \textbf{product} of the series.
\end{frame}

\begin{frame}{Examples}
$$\prod_{n=1}^\infty \left(1+\frac{1}{2^n}\right)$$
$$p_n = \left(1+\frac{1}{2}\right)\left(1+\frac{1}{4}\right)\left(1+\frac{1}{8}\right)\dots\left(1+\frac{1}{2^n}\right)$$
Can we compute $p_n$?
\end{frame}

\begin{frame}{Examples}\mbox{}
{\small
\begin{align*}
\left(1-\frac{1}{2}\right)p_n
  & = \left(1-\frac{1}{2}\right)\left(1+\frac{1}{2}\right)\left(1+\frac{1}{4}\right)\left(1+\frac{1}{8}\right)\dots\left(1+\frac{1}{2^n}\right)\\
  & = \left(1-\frac{1}{4}\right)\left(1+\frac{1}{4}\right)\left(1+\frac{1}{8}\right)\dots\left(1+\frac{1}{2^n}\right)\\
  & = \left(1-\frac{1}{8}\right)\left(1+\frac{1}{8}\right)\dots\left(1+\frac{1}{2^n}\right)\\
  & = 1-\frac{1}{2^n}\quad\Longrightarrow\quad p_n = 2-\frac{1}{2^{n-1}}
\end{align*}
}
$$\prod_{n=1}^\infty \left(1+\frac{1}{2^n}\right) = \lim_{n\rightarrow\infty} p_n = 2$$
\end{frame}

\begin{frame}{Examples}
Weierstrass product of sine
$$\sin(\pi x) = \pi z\prod_{n=1}^\infty \left(1-\frac{x^2}{n^2}\right)$$
\begin{itemize}
\item first thought of by Euler.
\item can be used to prove 
$$\sum \frac{1}{n^2} = \frac{\pi^2}{6}.$$
\end{itemize}
\end{frame}

\begin{frame}{Examples}
Let $q_1=2,q_2=3,q_3=5,\dots$ be the prime integers.\\
The infinite product
{\small
$$\zeta(s) = \left(\frac{1}{1-2^{-s}}\right)\left(\frac{1}{1-3^{-s}}\right)\left(\frac{1}{1-5^{-s}}\right)\left(\frac{1}{1-7^{-s}}\right)\left(\frac{1}{1-11^{-s}}\right)\dots$$
}
$$\zeta(s) = \prod_{n=1}^\infty \frac{1}{1-q_n^{-s}}$$
converges for $s>1$.\\
This product is called the \textbf{Riemann zeta function}.
\end{frame}

\begin{frame}{Riemann zeta}
There's an amazing way to convert to a series:
\begin{thm}
For any real number $s>1$
$$\zeta(s) = \prod_{n=1}^\infty \frac{1}{1-q_n^{-s}} = \sum_{m=1}^{\infty}\frac{1}{m^s}$$
\end{thm}
\end{frame}
\begin{frame}{Riemann zeta}
Here's why:
$$\frac{1}{2^s} \sum_{m=1}^{\infty}\frac{1}{m^s} = \sum_{m=1}^\infty \frac{1}{(2m)^s}$$
$$\left(1-\frac{1}{2^s}\right) \sum_{m=1}^{\infty}\frac{1}{m^s} = \sum_{2\nmid m}^\infty \frac{1}{m^s}$$
$$\left(1-\frac{1}{2^s}\right)\left(1-\frac{1}{3^s}\right)  \sum_{m=1}^{\infty}\frac{1}{m^s} = \sum_{2,3\nmid m}^\infty \frac{1}{m^s}$$
$$\left[\prod_{k=1}^n\left(1-\frac{1}{p_k^s}\right)\right] \sum_{m=1}^{\infty}\frac{1}{m^s} = \sum_{p_2,p_3,\dots, p_n\nmid m}^\infty \frac{1}{m^s}$$
\end{frame}
\begin{frame}{Riemann zeta}
Taking the limit as $n\rightarrow\infty$:
$$\zeta(s)^{-1}\sum_{m=1}^\infty \frac{1}{m^s} = \sum_{m\ \text{not divisible by any prime}} \frac{1}{m^s} = 1.$$
$$\zeta(s) = \sum_{m=1}^\infty \frac{1}{m^s}.$$
\end{frame}

\begin{frame}{Examples}
\begin{thm}
For any real number $s>1$,
$$\zeta(s) = \sum_{n=1}^\infty \frac{1}{n^s}.$$
\end{thm}
Therefore
$$\frac{1}{(1-1/2^2)(1-1/3^2)(1-1/5^2)(1-1/7^2)(1-1/11^2)\dots} = \frac{\pi^2}{6}.$$
\end{frame}

\begin{frame}{Relation to series}
\begin{thm}
Assume $a_n > 0$ for all $n$.
Then $\prod (1+a_n)$ converges if and only if $\sum a_n$ converges.
\end{thm}
\end{frame}

\begin{frame}{Relation to series}
\begin{defn}
We say that a product $\prod (1+a_n)$ \textbf{converges absolutely} if $\prod (1+|a_n|)$ converges.
\end{defn}
\begin{thm}
If $\prod (1+a_n)$ converges absolutely (equiv. $\sum a_n$ converges absolutely), then $\prod (1+a_n)$ converges.
\end{thm}
\end{frame}

\end{document}


